% !TEX program = lualatex
\documentclass[aspectratio=169,professionalfonts]{beamer}
\newcommand{\slidemode}{1}  % カラー投影モード
\usepackage{beamer_template}
\usepackage{enumerate}

\newcommand{\LectureNum}{6}
\newcommand{\LectureTitle}{たたみ込み}
\newcommand{\CourseName}{応用数学}
\renewcommand{\TermName}{前期}

\begin{document}

\begin{frame}{第6回　たたみ込み}
  \begin{block}{本回の内容}
  たたみ込み
  \end{block}
  \vfill\hfill{\scriptsize \textcolor{gray}{第2章 ラプラス変換　§2}}
\end{frame}

\begin{frame}{定義：たたみこみ（合成積）}
  \begin{block}{}
  \[
    (f * g)(t) = \int_{0}^t f(t-\tau)g(\tau) d\tau
  \]
  {\footnotesize \textcolor{gray}{$f(t), g(t)$ はいずれも $t \geq 0$ で定義された関数}}
  \end{block}
\end{frame}

\begin{frame}{定理：交換法則}
  \begin{block}{}
  \[
    f(t) * g(t) = g(t) * f(t)
  \]
  \end{block}
\end{frame}

\begin{frame}{定理：分配法則}
  \begin{block}{}
  \[
    f * (g_{1} + g_{2}) = f*g_{1} + f*g_{2}
  \]
  \end{block}
\end{frame}

\begin{frame}{定理：たたみこみのラプラス変換（たたみこみ定理）}
  \begin{block}{}
  \[
    L[f(t) * g(t)] = L[f(t)] \cdot L[g(t)] = F(s) \cdot G(s)
  \]
  \[
    L^{-1}[F(s) \cdot G(s)] = f(t) * g(t) = \int_{0}^t f(t-\tau)g(\tau) d\tau
  \]
  \end{block}
\end{frame}

\begin{frame}{例題5}
  \begin{exampleblock}{問題}
    $\int_{0}^t \sin \tau \cdot \cos(t-\tau) d\tau$ を求めよ\\[2pt]
  \end{exampleblock}
\end{frame}

\begin{frame}{例題5　解答}
  これは sin t * cos t のたたみこみ\\[4pt]
  積和公式 $: \sin \tau\cdot \cos(t-\tau) = (1/2)[\sin(\tau+(t-\tau)) + \sin(\tau-(t-\tau))]$\\[4pt]
  $= (1/2)[\sin t + \sin(2\tau-t)]$\\[4pt]
  $\int_{0}^t (1/2)[\sin t + \sin(2\tau-t)]d\tau$\\[4pt]
  $= (1/2)\begin{cases} [\sin t \cdot \tau]_{0}^t + [-\cos(2\tau-t)/2]_{0}^t \end{cases}$\\[4pt]
  $= (1/2)\begin{cases} t \sin t + & (-\cos t/2 - (-\cos(-t)/2)) \end{cases}$\\[4pt]
  $= (1/2){t \sin t + 0}$（$\cos(-t) = \cos t$ より）
  \medskip\textbf{答}：$\int_{0}^t \sin \tau \cdot \cos(t-\tau) d\tau = (t/2) \sin t$
\end{frame}

\begin{frame}{例題6}
  \begin{exampleblock}{問題}
    $f(t) * g(t) = g(t) * f(t)$ を証明せよ\\[2pt]
  \end{exampleblock}
\end{frame}

\begin{frame}{例題6　解答}
  $f(t)*g(t) = \int_{0}^t f(t-\tau)g(\tau)d\tau$\\[4pt]
  $u = t-\tau$ と置換 $: d\tau = -du, \tau=0 \to u=t, \tau=t \to u=0$\\[4pt]
  $= \int_{t}^0 f(u)g(t-u)(-du) = \int_{0}^t g(t-u)f(u)du$\\[4pt]
  $= g(t)*f(t)$
\end{frame}

\begin{frame}{例題7}
  \begin{exampleblock}{問題}
    $L[f(t)] = F(s)$ のとき、関数 $F(s)/(s-\alpha)$ の逆ラプラス変換を求めよ。ただし $\alpha$ は定数とする。\\[2pt]
  \end{exampleblock}
\end{frame}

\begin{frame}{例題7　解答}
  $L^{-1}[F(s)/(s-\alpha)] = L^{-1}[F(s) \cdot 1/(s-\alpha)]$\\[4pt]
  $L^{-1}[1/(s-\alpha)] = e^{\alpha t}$ より、たたみこみ定理を適用\\[4pt]
  $= L^{-1}[F(s)] * e^{\alpha t} = f(t) * e^{\alpha t}$
  \medskip\textbf{答}：$L^{-1}[F(s)/(s-\alpha)] = \int_{0}^t e^{\alpha(t-\tau)} f(\tau) d\tau$
\end{frame}

\begin{frame}{例題8}
  \begin{exampleblock}{問題}
    $\int_{0}^t x(\tau) \sin(t-\tau) d\tau = t^{2}$ を満たす関数 $x(t)$ を求めよ（積分方程式）\\[2pt]
  \end{exampleblock}
\end{frame}

\begin{frame}{例題8　解答}
  左辺は $x(t) * \sin t$ のたたみこみ\\[4pt]
  両辺をラプラス変換 $: X(s) \cdot L[\sin t] = L[t^{2}]$\\[4pt]
  $X(s) \cdot 1/(s^{2}+1) = 2/s^{3}$\\[4pt]
  $X(s) = 2(s^{2}+1)/s^{3} = 2/s + 2/s^{3}$\\[4pt]
  逆ラプラス変換 $: x(t) = 2 + t^{2}$
  \medskip\textbf{答}：$x(t) = t^{2} + 2$
\end{frame}

\begin{frame}{演習問題}
  \begin{exampleblock}{自分で解いてみよう}
  \begin{description}
    \item[問5] 次のたたみこみを求めよ
    $(1) \int_{0}^t e^\tau \sin(t-\tau) d\tau$\\
    $(2) \int_{0}^t \tau\cdot \cos(t-\tau) d\tau$\\
    \item[問6] 分配法則 $f*(g_{1}+g_{2}) = f*g_{1} + f*g_{2}$ を証明せよ
    \item[問7] $L[t^{2}*t]$ を求めよ。また、問 5 の結果から直接 $L[t^{2}*t]$ を求めよ
    \item[問8] $L[f(t)] = F(s)$ のとき、次の関数の逆ラプラス変換を求めよ。
    ただし $\alpha, \beta$ は定数で、 (2) では $\alpha \neq \beta$ 、 (3) では $\beta \neq 0$ とする。\\
    $(1) F(s)/(s-\alpha)^{2}$\\
    $(2) F(s)/((s-\alpha)(s-\beta))$\\
    $(3) F(s)/((s-\alpha)^{2}+\beta^{2})$\\
    \item[問9] 次の積分方程式を満たす関数 $x(t)$ を求めよ。
    $x(t) + \int_{0}^t x(t-\tau)e^\tau d\tau = \cos 3t$\\
  \end{description}
  \end{exampleblock}
\end{frame}

\end{document}
